\chapter{Grundlagen}
\label{kap:grundlagen}

Dieses Kapitel zeigt Abbildungen, Tabellen und Quellcode.

\section{Abbildungen und Tabellen}

Abbildungen werden zentriert, die Beschriftung steht darunter. Bei übernommenen
Abbildungen steht die Quelle in einer Fußnote am Ende der Beschriftung.

\begin{figure}[htb]
  \centering
  \includegraphics[width=0.5\textwidth]{example-image}
  \caption[Beispiel für eine Abbildung]{Beispiel für eine
    Abbildung\protect\footnotemark}
  \label{abb:beispiel}
\end{figure}
\footnotetext{Entnommen aus \textcite[17]{mustermann2023}.}

Bei Tabellen steht die Beschriftung darüber.

\begin{table}[htb]
  \caption{Beispiel für eine Tabelle}
  \label{tab:beispiel}
  \centering
  \begin{tabular}{lrr}
    \hline
    Kennzahl & 2025 & 2026 \\
    \hline
    Umsatz in Mio.~€ & 12,4 & 13,1 \\
    Mitarbeitende & 85 & 92 \\
    \hline
  \end{tabular}
\end{table}

\section{Quellcode}

Programmcode wird immer als Text gesetzt, nie als Bildschirmfoto.

\begin{lstlisting}[language=Python, caption={Beispiel für Quellcode}, label=lst:beispiel]
def umsatz_wachstum(alt: float, neu: float) -> float:
    """Wachstum in Prozent."""
    return (neu - alt) / alt * 100
\end{lstlisting}
